\section{Discussion}
\label{sec:discussion}

\paragraph{Graph completeness and tool synergy.}
The 31\,\% of baseline slots attributable to callback gaps illustrates the penalty
imposed when static call graphs omit function-pointer dispatches.
Under a standard 15-turn traversal, targets separated by callback boundaries remained
unreached.
However, our sensitivity analysis revealed that graph tools do not operate in a vacuum:
given an expanded 30-turn budget, \textsc{CT-TOOL} combined graph exploration with symbol
searches and source inspection to locate both Q1 and Q4.
Thus, static graph gaps represent a formidable barrier for pure edge traversals,
though hybrid agents equipped with source inspection can partially bridge them at
substantial token expense.

\paragraph{Compile-time vs.\ syntactic call graphs.}
A central empirical question is whether compiler-grade extraction (\clang) outperforms
lightweight syntactic extraction (tree-sitter).
Our results show no measurable difference in fault-localisation accuracy: both \textsc{CT}
and \textsc{SYN} conditions achieved identical exact accuracy across both models
(1 of 5 exact excluding Q2 on Qwen; 1 of 6 on Nemotron).
While \clang extracted 2{,}438 call-site edges with argument bindings compared to 772 unique
call pairs in \texttt{graphify}, the additional compiler precision did not translate into
higher localisation accuracy on this benchmark.

\paragraph{Text slices vs.\ interactive tool navigation.}
Text slices match or exceed tool conditions in exact accuracy at a fraction of
the token cost (3{,}400--5{,}500 tokens vs.\ 39{,}474--86{,}373 for tool conditions;
10--12$\times$ for Qwen, 16--23$\times$ for Nemotron).
In top-3 coverage, tool conditions are equal or strictly higher in three of four
comparisons, though the margin is only a single question.
For Q5 on \qwen, the BFS slice happens to include
\texttt{sqlite3\_str\_vappendf} at shallow depth, while the tool agent
stops at a callee one hop deeper (\texttt{printfTempBuf}).
This suggests that for shallow-target questions where the answer lies within a
small BFS radius of the entry, a text slice is a cost-effective alternative.

\paragraph{Caller/callee tools and edge citation validity.}
The not-in-graph rate was 2.4$\times$ lower in conditions with caller/callee
tools (12\,\%) than in source-only conditions (28\,\%).
We report this as a descriptive association on one model (133 edges across 36
runs); no statistical test was applied.
Some absent edges reflect function-pointer structural gaps rather than model
fabrication, so the rate is an upper bound on actual hallucination.
Even so, conditions with edge-traversal tools appear to constrain the model to
citing edges that actually exist in the graph, even when those tools do not
change the final answer.

\paragraph{Majority consensus and candidate comparison.}
Taking a simple majority consensus ($2/3$) across the three sensitivity repetitions
yields consistent majority predictions on two benchmark answers: Q1 on \textsc{CT-TOOL} (2 of 3 votes for
\texttt{resolveOrderByTermToExprList}) and Q5 on \textsc{SYN-TOOL} (2 of 3 votes for
\texttt{sqlite3\_str\_vappendf}).
While three samples do not establish statistical proof, this observation indicates that
sampling multiple trajectories with majority consensus can filter out transient misdirections
(such as naming \texttt{walkWindowList} or \texttt{printfTempBuf}).
Furthermore, qualitative inspection of failed trajectories suggests that prompting the
agent to explicitly compare two or three candidate functions before committing---rather
than selecting the first symptom-matching name such as \texttt{codeDistinct}---would
further reduce false positives.
